\chapter{El muro de \texorpdfstring{\ident{substfc}}{substfc}}
\label{cap:muro-substfc}

Los libros de incompletitud cuentan la demostración. Esta parte cuenta lo otro: qué pasa cuando
la pieza que falta no es difícil, sino que \textbf{la ruta obvia hacia ella hace inconsistente la
teoría}, y hay que descubrirlo antes de escribirla.

El episodio de este capítulo duró cuarenta días. No fue un atasco de implementación: fue un muro
con tres paredes independientes, y la primera de ellas —la que parecía la salida— habría destruido
el proyecto entero de haberse aceptado. Se atrapó porque alguien se sentó a derivar la
contradicción a mano en vez de comprobar que compilaba.

\section{Qué es \texorpdfstring{\ident{substfc}}{substfc}, y por qué está en todas partes}

El capítulo~\ref{cap:codigos} construyó los códigos de la sintaxis. Sustituir es lo que
la sintaxis hace: si $\varphi$ tiene una variable libre, $\varphi[t/v]$ es la fórmula que resulta de
poner $t$ donde estaba $v$. En la metateoría eso es la función \ident{substFormula}. Dentro de la
teoría objeto tiene que existir su contraparte, y se llama \ident{substfc}:

\leanfrag{substfc_def}

\selee{«\ident{substfc} $v$ $t$ $f$ es el término objeto \emph{«sustituir la variable de código $v$
por el término de código $t$ en la fórmula de código $f$»}»}

Nótese qué es y qué no es. No computa nada: construye un \emph{símbolo de función objeto} de aridad
tres. Lean no sabe qué significa. Quien lo sabe es la teoría, mediante ocho axiomas objeto, uno por
constructor de fórmula. Éste es el del condicional:

\leanfrag{ax_substfc_impl}

\selee{«para cualesquiera $v,t,a,b$: sustituir en el condicional $a \Rightarrow b$ es el condicional
de las dos sustituciones»}

Y aquí está la razón de que todo esto importe. El término diagonal —el corazón mismo de la sentencia
de Gödel, la que dice de sí misma que no es demostrable— \emph{es} una aplicación de
\ident{substfc}:

\leanfrag{diagTerm}

Si \ident{substfc} no funciona, no hay autorreferencia. Si tiene una grieta, la grieta está en el
centro del edificio.

\section{Lo que bloqueaba: siete de veintiuno}

D3 —la tercera condición de derivabilidad, la que Gödel anunció y nunca publicó, y sin la cual no
hay segundo teorema— se reduce en este proyecto a un solo lema, y ése a que el verificador de
pruebas se refleje dentro de $\Prov$: que la teoría demuestre, de cada línea de una derivación, que
la línea está bien formada.

Las líneas se clasifican por un \emph{tag}, y hay veintiún tags. Cada uno necesita su
\defterm{reflector}\nivelterm{proyecto}: el teorema que traslada al interior de $\Prov$ lo que la
metateoría ya sabe de ese tag. Catorce estaban hechos. Siete no.

Y aquí el proyecto hizo algo que conviene subrayar, porque es lo contrario de lo que se suele
hacer con una tarea pendiente: en vez de estimar cuánto faltaba, lo \textbf{demostró}.

\leanfrag{lineWF_modulo_7}

\begin{leccion}[Un teorema cuya función es medir lo que falta]
Ese enunciado no prueba nada nuevo sobre aritmética. Prueba que \emph{el ensamblaje encaja}: que
tomando los siete reflectores que faltan \textbf{como hipótesis}, todo lo demás cierra. Y por tanto
que cerrar el verificador es \textbf{exactamente} cerrar esos siete, ni una pieza más aguas abajo.

Es una manera de convertir «creo que falta poco» en un objeto que compila o no compila. Cuesta una
tarde y evita meses de trabajo en la dirección equivocada — que es justo lo que este capítulo
narra que estuvo a punto de pasar.
\end{leccion}

Los siete —\ident{q1}, \ident{q2}, \ident{q3}, \ident{leibniz}, \ident{ind}, \ident{qconf},
\ident{listInd}— tienen algo en común y sólo eso: su conclusión se reconstruye con \ident{substfc}
o con \ident{liftfc}, y por tanto exige \emph{evaluar} esos operadores dentro de $\Prov$. Eso es el
muro.

\section{El axioma que lo resolvía, y que hace inconsistente la teoría}

Hay una salida evidente, y es la primera que se le ocurre a cualquiera. El operador dotado
\ident{substfcT} es la versión «código del código» de \ident{substfc}. Bastaría postular que
conmutan:
\[
  \text{\ident{ax_tc_substfc}}:\quad
  \dotado{(\text{\ident{substfc}}\ v\ s\ f)} \;\doteq\;
  \text{\ident{substfcT}}\ \dotado{v}\ \dotado{s}\ \dotado{f}.
\]

\selee{«el código del resultado de sustituir es el resultado de sustituir sobre los códigos»}

Parece inocente. Ya existe un axioma exactamente de esa forma y nadie protesta — el de \ident{cons}:

\leanfrag{ax_tc_cons}

La diferencia es la que decide el episodio, y no se ve mirando los dos enunciados: \ident{cons} es
un \textbf{constructor}, y \ident{substfc} es una \textbf{función definida} — definida, además, por
ocho ecuaciones que la teoría \emph{ya tiene} y que la interpretan como productora de árboles de
\ident{cons}. Postular la ley de conmutación sobre algo ya interpretado no añade información:
la contradice.

\begin{refutado}[La derivación, en cinco pasos]
Supóngase \ident{ax_tc_substfc} en la teoría. Tómese la fórmula $\text{\ident{implc}}\ a\ b$.
\begin{enumerate}
\item Por \ident{ax_substfc_impl} —axioma ya presente—, la teoría demuestra que
  $\text{\ident{substfc}}\ v\ s\ (\text{\ident{implc}}\ a\ b)$ \emph{es} un
  $\text{\ident{implc}}$ de dos sustituciones. La teoría \textbf{interpreta} \ident{substfc}: no
  es opaco.
\item El operador $\dotado{\cdot}$ respeta la igualdad derivada —es un símbolo de función, y la
  congruencia de Leibniz vale para él—:
  \leanfrag{prf_congr_tcFn}
  Luego los códigos de los dos lados de (1) son iguales.
\item El lado izquierdo, por el axioma supuesto, tiene cabeza
  $\langle 1, \codigo{\text{\ident{substfc}}}, \ldots\rangle$.
\item El lado derecho es el código de un \ident{cons} —porque \ident{implc} \emph{es} un
  \ident{cons}—, y por \ident{ax_tc_cons} tiene cabeza
  $\langle 1, \codigo{\ident{::}}, \ldots\rangle$.
\item Luego la teoría demuestra que esas dos cabezas son iguales. Pero la aritmética negativa de
  códigos, que ya está construida, demuestra que \textbf{no} lo son:
  \leanfrag{cons_ne_head}
  Contradicción. La teoría prueba $\bot$.
\end{enumerate}
\end{refutado}

\begin{muro}[Por qué esto no habría fallado en compilación]
Una teoría objeto inconsistente \textbf{compila}. Lean no comprueba que los axiomas objeto sean
satisfacibles: son datos, valores de tipo \ident{Formula}, y añadir uno más a la lista es tan legal
como añadir el primero. El proyecto habría seguido en verde durante meses, demostrando teoremas
—todos ellos vacíos, porque de una teoría inconsistente se demuestra cualquier cosa—. Y en
particular habría «demostrado» $\Prov(\codigo{\bot})$, que aniquila la construcción de Gödel entera:
la sentencia $G$ dejaría de decir nada.

Este es el modo de fallo característico del proyecto, y por eso la Parte~\ref{parte:noescrito}
existe. Ya apareció en pequeño en el capítulo~\ref{cap:kernel} —un axioma del kernel que era falso—
y volverá a aparecer en grande en el capítulo~\ref{cap:numerales}, donde no se atrapó a tiempo.
\end{muro}

\begin{observacion}
\ident{ax_tc_substfc} \textbf{no existe} en el código. Y sin embargo su nombre aparece
\textbf{nueve veces} en el árbol —una en \modulo{ROBINSON_PlusPlus/Meta/EvalArithPrf.lean}, otra en
\modulo{ROBINSON_PlusPlus/Meta/SubstfcCodePrf.lean}, el resto en sondeos y experimentos—: siempre
en un comentario, siempre para advertir de que sería inconsistente. El proyecto convirtió la
refutación en un cartel y lo clavó en las nueve puertas por las que alguien podría volver a entrar.
Es, en pequeño, la misma doctrina que este libro: un resultado negativo también es un resultado, y
se guarda donde vaya a leerse.
\end{observacion}

\section{El obstáculo real: evaluar sobre un argumento abstracto}

Cerrada la puerta del axioma, queda el trabajo honesto: demostrar la evaluación en vez de
postularla. Y ahí está la segunda pared.

Hay evaluaciones provables que salen. La de \ident{carc} —la cabeza de una lista— sale porque
\ident{carc} tiene \textbf{una} ecuación definitoria y se instancia sobre un patrón explícito:
delante hay un \ident{cons} escrito, y se aplica la ecuación. \ident{substfc} no está en esa
situación por dos razones a la vez: se define por recursión sobre los \textbf{ocho} constructores
de fórmula, con ligadores —los casos $\forall$ y $\exists$ incrementan el nivel y desplazan el
sustituyendo—, y en los siete tags se aplica a un código \textbf{abstracto}, una casilla de la
línea que se está verificando, no a una fórmula escrita.

\begin{muro}[Lo que eso obliga a hacer]
Evaluar \ident{substfc} sobre un argumento abstracto exige \textbf{inducción sobre códigos de
fórmula dentro de $\Prov$}. No inducción en la metateoría —eso se tiene— sino inducción de la
teoría objeto sobre sus propios códigos, disponible en un punto en que la teoría objeto no tiene
inducción: la aritmética de \modulo{ROBINSON_PlusPlus/Minimal/} es Q\texttt{++}, sin esquema de
inducción, por la razón que dio el capítulo~\ref{cap:qpp}.

Ésa es la distancia real entre «un lema más» y «un problema abierto», y explica por qué el registro
del proyecto insiste en que el muro no era mecánico.
\end{muro}

\section{La tercera pared: quitar \texorpdfstring{\ident{substfc}}{substfc} tampoco se puede}

Queda una salida más, y es la que parece más razonable de todas: si el problema es que los siete
esquemas mencionan \ident{substfc}, reescríbanse sin él. El 25 de julio de 2026 se sondeó
sistemáticamente, y el resultado no fue «no lo hemos conseguido» sino una dicotomía cerrada.

El argumento tiene dos bisagras, las dos comprobables en el código: la regla de instanciación toma
un término \textbf{arbitrario}, y \ident{substFormula} es recursión sobre los ocho constructores.
De ahí: para que el reflector sea \emph{sólido} tiene que atar la instancia correcta, y la instancia
correcta contiene $\varphi[t/v]$. Si esa mención queda \textbf{ligada} —como función, como ecuación,
como átomo relacional—, reaparece el mismo muro un piso más abajo. Si queda \textbf{libre}, el
esquema deja de ser sólido. Las tres propuestas concretas que se estudiaron colapsan las tres al
mismo núcleo.

\begin{leccion}[La diferencia entre «no sale» y «no existe»]
Un resultado como éste es incómodo de publicar y valiosísimo de tener: no dice que el equipo no
supiera hacerlo, dice que no hay nada que hacer por ahí. Sin él, la ruta se habría reintentado cada
pocas semanas con una variante — que es exactamente lo que pasa cuando un callejón sin salida no se
documenta como tal. Los otros casos de esta misma clase que el proyecto encontró después los
recoge \capfuturo{las tareas que resultan ser imposibles}.
\end{leccion}

\section{Cuarenta días}

\begin{center}\small
\begin{tabular}{lll}
\toprule
fecha & commit & hito \\
\midrule
2026-07-10 & \texttt{8a96e8e} & el frente arranca en positivo \\
\textbf{2026-07-22} & \texttt{15dccda} & \textbf{el muro se identifica}: «el plan era erróneo en su forma» \\
2026-07-24 & --- & veredicto: \ident{ax_tc_substfc} es inconsistente; no hay ruta \defterm{net-0}\nivelterm{proyecto} \\
2026-07-25 & \texttt{6c4d7ac} & la puerta de reescribir los esquemas, cerrada \\
2026-07-26 & \texttt{ee31d07} & inducción fuerte objeto, net-0 — la herramienta que luego serviría \\
2026-08-25 & \texttt{bdfff12} & el reconocedor fusionado no desbloquea el muro \\
2026-08-30 & \texttt{2f27a29} & el descenso, cerrado \\
\textbf{2026-08-31} & \texttt{c1fd804} & \textbf{el muro está roto} \\
2026-09-07 & \texttt{61edd46} & dentro del build: deja de ser un sondeo y pasa a ser un teorema \\
\bottomrule
\end{tabular}
\end{center}

Cuarenta días entre la identificación y la rotura; siete más hasta que el resultado dejó de vivir
en \modulo{sondeos/} y entró en el build, que es la condición que este libro exige para llamarlo
teorema (§2.2). Hoy se enuncia así:

\leanfrag{pcc_eval_substfc}

\selee{«bajo dos guardas —que $s$ tenga testigo y que $f$ sea un código de fórmula bien formado
contra su testigo—, la teoría demuestra que el operador dotado aplicado a los códigos es igual al
código del resultado»}

Las dos hipótesis no son deudas colgando: son las \emph{guardas} que dicen sobre qué códigos vale
la evaluación, y \capfuturo{la no-vacuidad de las guardas} muestra que se satisfacen
para códigos reales. Cómo se llegó hasta aquí —definir en vez de axiomatizar, partir el
reconocedor en tres, y una inducción con la guarda dentro— es la materia del resto de esta
parte.

\begin{observacion}
Un detalle que este libro no puede pasar por alto, porque es su propia doctrina aplicada a sí mismo.
El docstring de \ident{pcc_eval_substfc} dice, literalmente, «cero \ident{sorry}, cero hipótesis».
El teorema tiene dos. Es una imprecisión menor —«hipótesis» quiere decir ahí «obligación
pendiente»—, pero es exactamente el tipo de afirmación que un lector tomaría por buena y que el
compilador no verifica. Queda anotada en \doc{DOCSTRINGS-NO-FIABLES.md}.
\end{observacion}
